% ============================================================================
% BỔ SUNG ĐỢT 2 — TOÁN 9 KNTT CHƯƠNG V: BÀI 16 & BÀI 17 (30 CÂU)
% Mỗi mục nhỏ thêm 2 TN Nhận biết + 2 Đúng/Sai Thông hiểu + 2 TLN Vận dụng
% ============================================================================

% --- 9H5?16-1: Xác định vị trí tương đối của đường thẳng và đường tròn ---

\begin{ex}%[9H5N16-1]
% nguon: SGK Toán 9 KNTT tập 1 – Bài 16, Vị trí tương đối của đường thẳng và đường tròn
	Nếu đường thẳng $d$ cắt đường tròn $(O; R)$ thì khoảng cách $h$ từ tâm $O$ đến đường thẳng $d$ thỏa mãn
	\choice
	{\True $h < R$}
	{$h = R$}
	{$h > R$}
	{$h \ge 2R$}
	\loigiai{
		Đường thẳng $d$ cắt đường tròn $(O; R)$ khi và chỉ khi khoảng cách từ tâm $O$ đến $d$ nhỏ hơn bán kính, tức là $h < R$.
	}
\end{ex}

\begin{ex}%[9H5N16-1]
% nguon: SGK Toán 9 KNTT tập 1 – Bài 16, Vị trí tương đối của đường thẳng và đường tròn
	Cho đường tròn $(O; 8\text{ cm})$ và đường thẳng $a$ có khoảng cách từ tâm $O$ đến $a$ bằng $8\text{ cm}$. Số điểm chung của đường thẳng $a$ và đường tròn $(O)$ là
	\choice
	{$0$}
	{\True $1$}
	{$2$}
	{Vô số}
	\loigiai{
		Vì khoảng cách từ tâm $O$ đến đường thẳng $a$ bằng bán kính $R = 8\text{ cm}$ nên $a$ tiếp xúc với $(O)$, do đó chúng có đúng $1$ điểm chung.
	}
\end{ex}

\begin{ex}%[9H5H16-1]
% nguon: SGK Toán 9 KNTT tập 1 – Bài 16, Vị trí tương đối của đường thẳng và đường tròn
	Trong mặt phẳng tọa độ $Oxy$, cho đường tròn tâm $M(4; 3)$ bán kính $R = 4$. Xét tính đúng, sai của các khẳng định sau:
	\choiceTF
	{\True Khoảng cách từ tâm $M$ đến trục hoành $Ox$ bằng $3$}
	{\True Khoảng cách từ tâm $M$ đến trục tung $Oy$ bằng $4$}
	{Đường tròn $(M; 4)$ không giao nhau với trục hoành $Ox$}
	{\True Đường tròn $(M; 4)$ tiếp xúc với trục tung $Oy$ và cắt trục hoành $Ox$ theo một dây cung có độ dài bằng $2\sqrt{7}$}
	\loigiai{
		\begin{itemchoice}
			\itemch Đúng. Khoảng cách từ $M(4; 3)$ đến trục hoành $Ox$ là $|y_M| = 3$.
			\itemch Đúng. Khoảng cách từ $M(4; 3)$ đến trục tung $Oy$ là $|x_M| = 4$.
			\itemch Sai. Vì khoảng cách từ $M$ đến $Ox$ là $3 < R = 4$ nên đường tròn $(M; 4)$ cắt trục hoành $Ox$ tại hai điểm phân biệt.
			\itemch Đúng. Vì $d(M, Oy) = 4 = R$ nên $(M; 4)$ tiếp xúc với $Oy$; độ dài dây cung cắt trên $Ox$ là $2\sqrt{R^2 - d^2} = 2\sqrt{4^2 - 3^2} = 2\sqrt{7}$.
		\end{itemchoice}
	}
\end{ex}

\begin{ex}%[9H5H16-1]
% nguon: SGK Toán 9 KNTT tập 1 – Bài 16, Vị trí tương đối của đường thẳng và đường tròn
	Cho đường tròn $(O; 5\text{ cm})$ và đường thẳng $d$ cách tâm $O$ một khoảng $h = 3\text{ cm}$. Xét tính đúng, sai của các khẳng định sau:
	\choiceTF
	{\True Đường thẳng $d$ cắt đường tròn $(O)$ tại hai điểm phân biệt $A, B$}
	{\True Độ dài dây cung $AB$ bằng $8\text{ cm}$}
	{Nếu dịch chuyển đường thẳng $d$ ra xa tâm $O$ thêm $1\text{ cm}$ thì đường thẳng mới tiếp xúc với $(O)$}
	{\True Diện tích tam giác $OAB$ bằng $12\text{ cm}^2$}
	\loigiai{
		\begin{itemchoice}
			\itemch Đúng. Vì $h = 3 < R = 5$ nên $d$ cắt $(O)$ tại hai điểm phân biệt $A, B$.
			\itemch Đúng. Độ dài dây $AB = 2\sqrt{5^2 - 3^2} = 2 \cdot 4 = 8\text{ cm}$.
			\itemch Sai. Khi dịch chuyển ra xa tâm thêm $1\text{ cm}$ thì khoảng cách mới là $3 + 1 = 4\text{ cm} < 5\text{ cm}$, đường thẳng vẫn cắt đường tròn.
			\itemch Đúng. Diện tích tam giác $OAB$ là $\dfrac{1}{2} \cdot AB \cdot h = \dfrac{1}{2} \cdot 8 \cdot 3 = 12\text{ cm}^2$.
		\end{itemchoice}
	}
\end{ex}

\begin{bt}%[9H5V16-1]
% nguon: SGK Toán 9 KNTT tập 1 – Bài 16, Vị trí tương đối của đường thẳng và đường tròn
	Trong mặt phẳng tọa độ $Oxy$, cho điểm $B(-5; 6)$. Gọi $R$ là bán kính của đường tròn tâm $B$ sao cho đường tròn đó tiếp xúc với trục hoành $Ox$. Giá trị của $R$ bằng bao nhiêu?
	\shortans{6}
	\loigiai{
		Khoảng cách từ điểm $B(-5; 6)$ đến trục hoành $Ox$ bằng $|y_B| = |6| = 6$.\\
		Để đường tròn $(B; R)$ tiếp xúc với trục hoành $Ox$ thì $R = d(B, Ox) = 6$.
	}
\end{bt}

\begin{bt}%[9H5V16-1]
% nguon: SGK Toán 9 KNTT tập 1 – Bài 16, Vị trí tương đối của đường thẳng và đường tròn
	Cho đường tròn $(O; 10\text{ cm})$ và đường thẳng $d$ cách tâm $O$ một khoảng bằng $6\text{ cm}$. Đường thẳng $d$ cắt đường tròn $(O)$ tại hai điểm $A$ và $B$. Tính diện tích tam giác $OAB$ (đơn vị $\text{cm}^2$).
	\shortans{48}
	\loigiai{
		Gọi $H$ là hình chiếu của $O$ trên $d$, ta có $OH = 6\text{ cm}$ và $AH = \sqrt{OA^2 - OH^2} = \sqrt{10^2 - 6^2} = 8\text{ cm}$.\\
		Suy ra $AB = 2AH = 16\text{ cm}$.\\
		Diện tích tam giác $OAB$ là $S = \dfrac{1}{2} \cdot AB \cdot OH = \dfrac{1}{2} \cdot 16 \cdot 6 = 48\text{ cm}^2$.
	}
\end{bt}

% --- 9H5?16-2: Tiếp tuyến: dấu hiệu nhận biết, chứng minh tiếp tuyến ---

\begin{ex}%[9H5N16-2]
% nguon: SGK Toán 9 KNTT tập 1 – Bài 16, Vị trí tương đối của đường thẳng và đường tròn
	Nếu đường thẳng $a$ là tiếp tuyến của đường tròn $(O)$ tại điểm $M$ thì khẳng định nào sau đây đúng?
	\choice
	{$a$ đi qua tâm $O$}
	{\True $a \perp OM$ tại $M$}
	{$a$ song song với $OM$}
	{$OM$ là đường kính của $(O)$}
	\loigiai{
		Theo tính chất tiếp tuyến của đường tròn, nếu đường thẳng $a$ tiếp xúc với đường tròn $(O)$ tại $M$ thì $a$ vuông góc với bán kính $OM$ tại tiếp điểm $M$.
	}
\end{ex}

\begin{ex}%[9H5N16-2]
% nguon: SGK Toán 9 KNTT tập 1 – Bài 16, Vị trí tương đối của đường thẳng và đường tròn
	Qua một điểm $A$ nằm trên đường tròn $(O)$, ta có thể kẻ được bao nhiêu tiếp tuyến với đường tròn $(O)$?
	\choice
	{$0$}
	{\True $1$}
	{$2$}
	{Vô số}
	\loigiai{
		Qua một điểm $A$ nằm trên đường tròn $(O)$ chỉ có duy nhất một đường thẳng vuông góc với bán kính $OA$ tại $A$, nên chỉ kẻ được đúng $1$ tiếp tuyến.
	}
\end{ex}

\begin{ex}%[9H5H16-2]
% nguon: SGK Toán 9 KNTT tập 1 – Bài 16, Vị trí tương đối của đường thẳng và đường tròn
	Cho đường tròn $(O; 6\text{ cm})$ và điểm $M$ nằm ngoài đường tròn sao cho $OM = 10\text{ cm}$. Kẻ tiếp tuyến $MA$ đến $(O)$ với $A$ là tiếp điểm. Xét tính đúng, sai của các khẳng định sau:
	\choiceTF
	{\True Tam giác $OAM$ vuông tại $A$}
	{\True Độ dài đoạn thẳng $MA$ bằng $8\text{ cm}$}
	{Diện tích tam giác $OAM$ bằng $48\text{ cm}^2$}
	{\True Đường cao $AH$ kẻ từ $A$ xuống $OM$ có độ dài bằng $4{,}8\text{ cm}$}
	\loigiai{
		\begin{itemchoice}
			\itemch Đúng. Vì $MA$ là tiếp tuyến của $(O)$ tại $A$ nên $MA \perp OA$, do đó tam giác $OAM$ vuông tại $A$.
			\itemch Đúng. Áp dụng định lí Pythagore: $MA = \sqrt{OM^2 - OA^2} = \sqrt{10^2 - 6^2} = 8\text{ cm}$.
			\itemch Sai. Diện tích tam giác vuông $OAM$ là $\dfrac{1}{2} \cdot OA \cdot MA = \dfrac{1}{2} \cdot 6 \cdot 8 = 24\text{ cm}^2$.
			\itemch Đúng. Ta có $AH \cdot OM = OA \cdot MA \Rightarrow AH = \dfrac{6 \cdot 8}{10} = 4{,}8\text{ cm}$.
		\end{itemchoice}
	}
\end{ex}

\begin{ex}%[9H5H16-2]
% nguon: SGK Toán 9 KNTT tập 1 – Bài 16, Vị trí tương đối của đường thẳng và đường tròn
	Cho tam giác $ABC$ vuông tại $A$ có $AB = 5\text{ cm}$, $AC = 12\text{ cm}$. Xét đường tròn tâm $B$ bán kính $R = 5\text{ cm}$. Xét tính đúng, sai của các khẳng định sau:
	\choiceTF
	{\True Điểm $A$ nằm trên đường tròn $(B; 5\text{ cm})$}
	{\True Đường thẳng $AC$ là tiếp tuyến của đường tròn $(B; 5\text{ cm})$}
	{Độ dài cạnh $BC$ bằng $15\text{ cm}$}
	{\True Điểm $C$ nằm ngoài đường tròn $(B; 5\text{ cm})$ và khoảng cách từ $C$ đến điểm gần nhất trên đường tròn $(B; 5\text{ cm})$ bằng $8\text{ cm}$}
	\loigiai{
		\begin{itemchoice}
			\itemch Đúng. Vì $BA = 5\text{ cm} = R$ nên $A$ thuộc đường tròn $(B; 5\text{ cm})$.
			\itemch Đúng. Vì $A \in (B; 5\text{ cm})$ và $AC \perp BA$ tại $A$ nên $AC$ là tiếp tuyến của $(B; 5\text{ cm})$.
			\itemch Sai. Trong tam giác $ABC$ vuông tại $A$: $BC = \sqrt{5^2 + 12^2} = 13\text{ cm}$.
			\itemch Đúng. Vì $BC = 13 > 5$ nên $C$ nằm ngoài $(B; 5\text{ cm})$, khoảng cách ngắn nhất từ $C$ đến đường tròn là $BC - R = 13 - 5 = 8\text{ cm}$.
		\end{itemchoice}
	}
\end{ex}

\begin{bt}%[9H5V16-2]
% nguon: SGK Toán 9 KNTT tập 1 – Bài 16, Vị trí tương đối của đường thẳng và đường tròn
	Cho đường tròn $(O; 8\text{ cm})$. Từ điểm $P$ nằm ngoài đường tròn kẻ tiếp tuyến $PA$ ($A$ là tiếp điểm) sao cho $PA = 15\text{ cm}$. Tính độ dài đoạn thẳng $OP$ (đơn vị cm).
	\shortans{17}
	\loigiai{
		Vì $PA$ là tiếp tuyến của $(O)$ tại $A$ nên tam giác $OAP$ vuông tại $A$.\\
		Theo định lí Pythagore: $OP = \sqrt{OA^2 + PA^2} = \sqrt{8^2 + 15^2} = \sqrt{289} = 17\text{ cm}$.
	}
\end{bt}

\begin{bt}%[9H5V16-2]
% nguon: SGK Toán 9 KNTT tập 1 – Bài 16, Vị trí tương đối của đường thẳng và đường tròn
	Từ điểm $M$ nằm ngoài đường tròn $(O; 9\text{ cm})$ sao cho $OM = 15\text{ cm}$, kẻ tiếp tuyến $MA$ đến đường tròn $(O)$ với $A$ là tiếp điểm. Tính diện tích tam giác $OAM$ (đơn vị $\text{cm}^2$).
	\shortans{54}
	\loigiai{
		Tam giác $OAM$ vuông tại $A$ nên $MA = \sqrt{OM^2 - OA^2} = \sqrt{15^2 - 9^2} = 12\text{ cm}$.\\
		Diện tích tam giác $OAM$ là $S = \dfrac{1}{2} \cdot OA \cdot MA = \dfrac{1}{2} \cdot 9 \cdot 12 = 54\text{ cm}^2$.
	}
\end{bt}

% --- 9H5?16-3: Tính chất hai tiếp tuyến cắt nhau ---

\begin{ex}%[9H5N16-3]
% nguon: SGK Toán 9 KNTT tập 1 – Bài 16, Vị trí tương đối của đường thẳng và đường tròn
	Từ điểm $A$ nằm ngoài đường tròn $(O)$, kẻ hai tiếp tuyến $AB, AC$ đến đường tròn ($B, C$ là các tiếp điểm). Nếu $AB = 7\text{ cm}$ thì độ dài $AC$ bằng
	\choice
	{$3{,}5\text{ cm}$}
	{\True $7\text{ cm}$}
	{$14\text{ cm}$}
	{$7\sqrt{2}\text{ cm}$}
	\loigiai{
		Theo tính chất hai tiếp tuyến cắt nhau, điểm $A$ cách đều hai tiếp điểm $B$ và $C$, nên $AC = AB = 7\text{ cm}$.
	}
\end{ex}

\begin{ex}%[9H5N16-3]
% nguon: SGK Toán 9 KNTT tập 1 – Bài 16, Vị trí tương đối của đường thẳng và đường tròn
	Từ điểm $M$ nằm ngoài đường tròn $(O)$, kẻ hai tiếp tuyến $MA, MB$ đến đường tròn ($A, B$ là các tiếp điểm). Biết $\widehat{AMO} = 25^\circ$, số đo góc $\widehat{AMB}$ bằng
	\choice
	{$25^\circ$}
	{\True $50^\circ$}
	{$65^\circ$}
	{$130^\circ$}
	\loigiai{
		Theo tính chất hai tiếp tuyến cắt nhau, $MO$ là tia phân giác của $\widehat{AMB}$, suy ra $\widehat{AMB} = 2\widehat{AMO} = 2 \cdot 25^\circ = 50^\circ$.
	}
\end{ex}

\begin{ex}%[9H5H16-3]
% nguon: SGK Toán 9 KNTT tập 1 – Bài 16, Vị trí tương đối của đường thẳng và đường tròn
	Từ điểm $A$ nằm ngoài đường tròn $(O; 5\text{ cm})$ sao cho $OA = 13\text{ cm}$, kẻ hai tiếp tuyến $AB, AC$ đến $(O)$ với $B, C$ là các tiếp điểm. Xét tính đúng, sai của các khẳng định sau:
	\choiceTF
	{\True $AB = AC = 12\text{ cm}$}
	{\True $OA$ là đường trung trực của đoạn thẳng $BC$}
	{Diện tích tứ giác $OBAC$ bằng $120\text{ cm}^2$}
	{\True Chu vi tứ giác $OBAC$ bằng $34\text{ cm}$}
	\loigiai{
		\begin{itemchoice}
			\itemch Đúng. Trong tam giác $OBA$ vuông tại $B$: $AB = \sqrt{13^2 - 5^2} = 12\text{ cm}$, suy ra $AC = AB = 12\text{ cm}$.
			\itemch Đúng. Vì $OB = OC$ và $AB = AC$ nên $OA$ là đường trung trực của $BC$.
			\itemch Sai. Diện tích tứ giác $OBAC$ bằng $2S_{OBA} = 2 \cdot \dfrac{1}{2} \cdot 5 \cdot 12 = 60\text{ cm}^2$.
			\itemch Đúng. Chu vi tứ giác $OBAC$ là $OB + OC + AB + AC = 5 + 5 + 12 + 12 = 34\text{ cm}$.
		\end{itemchoice}
	}
\end{ex}

\begin{ex}%[9H5H16-3]
% nguon: SGK Toán 9 KNTT tập 1 – Bài 16, Vị trí tương đối của đường thẳng và đường tròn
	Từ điểm $M$ nằm ngoài đường tròn $(O; 4\text{ cm})$, kẻ hai tiếp tuyến $MA, MB$ vuông góc với nhau ($A, B$ là các tiếp điểm). Xét tính đúng, sai của các khẳng định sau:
	\choiceTF
	{\True Tứ giác $OAMB$ là hình vuông}
	{\True Độ dài $MA = MB = 4\text{ cm}$}
	{Độ dài đoạn thẳng $OM$ bằng $8\text{ cm}$}
	{\True Diện tích phần nằm ngoài đường tròn $(O)$ nhưng bên trong tứ giác $OAMB$ (lấy $\pi \approx 3{,}14$, làm tròn đến hàng phần trăm) bằng $3{,}44\text{ cm}^2$}
	\loigiai{
		\begin{itemchoice}
			\itemch Đúng. Tứ giác $OAMB$ có $\widehat{A} = \widehat{B} = \widehat{M} = 90^\circ$ và $OA = OB = 4\text{ cm}$ nên là hình vuông cạnh $4\text{ cm}$.
			\itemch Đúng. Vì $OAMB$ là hình vuông cạnh $4\text{ cm}$ nên $MA = MB = 4\text{ cm}$.
			\itemch Sai. Đường chéo hình vuông cạnh $4\text{ cm}$ là $OM = 4\sqrt{2}\text{ cm}$.
			\itemch Đúng. Phần diện tích cần tìm bằng hiệu diện tích hình vuông $OAMB$ và hình quạt tròn $OAB$ góc $90^\circ$: $4^2 - \dfrac{90}{360}\pi \cdot 4^2 = 16 - 4\pi \approx 16 - 4 \cdot 3{,}14 = 3{,}44\text{ cm}^2$.
		\end{itemchoice}
	}
\end{ex}

\begin{bt}%[9H5V16-3]
% nguon: SGK Toán 9 KNTT tập 1 – Bài 16, Vị trí tương đối của đường thẳng và đường tròn
	Từ điểm $M$ nằm ngoài đường tròn $(O; 3\text{ cm})$ sao cho $OM = 5\text{ cm}$, kẻ hai tiếp tuyến $MA, MB$ đến đường tròn $(O)$ ($A, B$ là các tiếp điểm). Tính diện tích tứ giác $OAMB$ (đơn vị $\text{cm}^2$).
	\shortans{12}
	\loigiai{
		Trong tam giác $OAM$ vuông tại $A$: $MA = \sqrt{5^2 - 3^2} = 4\text{ cm}$.\\
		Vì $\triangle OAM = \triangle OBM$ nên diện tích tứ giác $OAMB$ là:
		$$S_{OAMB} = 2S_{OAM} = 2 \cdot \dfrac{1}{2} \cdot OA \cdot MA = 3 \cdot 4 = 12\text{ cm}^2.$$
	}
\end{bt}

\begin{bt}%[9H5V16-3]
% nguon: SGK Toán 9 KNTT tập 1 – Bài 16, Vị trí tương đối của đường thẳng và đường tròn
	Cho đường tròn $(O)$ nội tiếp tam giác $ABC$, tiếp xúc với các cạnh $BC, CA, AB$ lần lượt tại các điểm $D, E, F$. Biết $AB = 7\text{ cm}$, $BC = 9\text{ cm}$, $CA = 8\text{ cm}$. Tính độ dài đoạn thẳng $AE$ (đơn vị cm).
	\shortans{3}
	\loigiai{
		Theo tính chất hai tiếp tuyến cắt nhau, đặt $AE = AF = x$, $BF = BD = y$, $CD = CE = z$.\\
		Ta có $x + y = AB = 7$, $y + z = BC = 9$, $z + x = CA = 8$.\\
		Cộng vế theo vế được $2(x + y + z) = 24 \Rightarrow x + y + z = 12$.\\
		Suy ra $AE = x = 12 - (y + z) = 12 - 9 = 3\text{ cm}$.
	}
\end{bt}

% --- 9H5?17-1: Xác định vị trí tương đối của hai đường tròn ---

\begin{ex}%[9H5N17-1]
% nguon: SGK Toán 9 KNTT tập 1 – Bài 17, Vị trí tương đối của hai đường tròn
	Cho hai đường tròn $(O; 9\text{ cm})$ và $(O'; 4\text{ cm})$ có khoảng cách giữa hai tâm là $OO' = 5\text{ cm}$. Vị trí tương đối của hai đường tròn đó là
	\choice
	{Ở ngoài nhau}
	{Tiếp xúc ngoài}
	{Cắt nhau}
	{\True Tiếp xúc trong}
	\loigiai{
		Ta có $R - r = 9 - 4 = 5\text{ cm} = OO'$, do đó hai đường tròn tiếp xúc trong với nhau.
	}
\end{ex}

\begin{ex}%[9H5N17-1]
% nguon: SGK Toán 9 KNTT tập 1 – Bài 17, Vị trí tương đối của hai đường tròn
	Cho hai đường tròn $(O; 6\text{ cm})$ và $(O'; 4\text{ cm})$ có khoảng cách giữa hai tâm là $OO' = 12\text{ cm}$. Vị trí tương đối của hai đường tròn đó là
	\choice
	{\True Ở ngoài nhau}
	{Tiếp xúc ngoài}
	{Cắt nhau}
	{Đựng nhau}
	\loigiai{
		Ta có $OO' = 12\text{ cm} > R + r = 6 + 4 = 10\text{ cm}$, nên hai đường tròn ở ngoài nhau.
	}
\end{ex}

\begin{ex}%[9H5H17-1]
% nguon: SGK Toán 9 KNTT tập 1 – Bài 17, Vị trí tương đối của hai đường tròn
	Cho hai đường tròn $(O; 10\text{ cm})$ và $(O'; 6\text{ cm})$ có khoảng cách giữa hai tâm là $OO' = d$. Xét tính đúng, sai của các khẳng định sau:
	\choiceTF
	{\True Hai đường tròn tiếp xúc ngoài khi $d = 16\text{ cm}$}
	{\True Hai đường tròn tiếp xúc trong khi $d = 4\text{ cm}$}
	{Nếu $d = 10\text{ cm}$ thì hai đường tròn ở ngoài nhau}
	{\True Có đúng $11$ giá trị nguyên của $d$ (tính theo cm) để hai đường tròn cắt nhau}
	\loigiai{
		\begin{itemchoice}
			\itemch Đúng. Hai đường tròn tiếp xúc ngoài khi $d = R + r = 10 + 6 = 16\text{ cm}$.
			\itemch Đúng. Hai đường tròn tiếp xúc trong khi $d = R - r = 10 - 6 = 4\text{ cm}$.
			\itemch Sai. Vì $4 < 10 < 16$ nên khi $d = 10\text{ cm}$ hai đường tròn cắt nhau.
			\itemch Đúng. Hai đường tròn cắt nhau khi $4 < d < 16$; các giá trị nguyên của $d$ là $5, 6, \ldots, 15$, gồm $15 - 5 + 1 = 11$ giá trị.
		\end{itemchoice}
	}
\end{ex}

\begin{ex}%[9H5H17-1]
% nguon: SGK Toán 9 KNTT tập 1 – Bài 17, Vị trí tương đối của hai đường tròn
	Cho hai đường tròn $(O; 5\text{ cm})$ và $(O'; 5\text{ cm})$ có khoảng cách hai tâm $OO' = 6\text{ cm}$ cắt nhau tại hai điểm $A$ và $B$. Xét tính đúng, sai của các khẳng định sau:
	\choiceTF
	{\True Hai đường tròn cắt nhau tại hai điểm phân biệt}
	{\True Đường nối tâm $OO'$ là đường trung trực của dây chung $AB$}
	{Độ dài dây chung $AB$ bằng $6\text{ cm}$}
	{\True Tứ giác $OAO'B$ là hình thoi có diện tích bằng $24\text{ cm}^2$}
	\loigiai{
		\begin{itemchoice}
			\itemch Đúng. Ta có $5 - 5 = 0 < 6 < 5 + 5 = 10$ nên hai đường tròn cắt nhau tại hai điểm phân biệt.
			\itemch Đúng. Vì $OA = OB = 5\text{ cm}$ và $O'A = O'B = 5\text{ cm}$ nên $OO'$ là đường trung trực của $AB$.
			\itemch Sai. Gọi $H = OO' \cap AB$, ta có $OH = \dfrac{OO'}{2} = 3\text{ cm}$ nên $AH = \sqrt{5^2 - 3^2} = 4\text{ cm}$, suy ra $AB = 8\text{ cm}$.
			\itemch Đúng. Vì $OA = AO' = O'B = BO = 5\text{ cm}$ nên $OAO'B$ là hình thoi, có diện tích $\dfrac{1}{2} \cdot OO' \cdot AB = \dfrac{1}{2} \cdot 6 \cdot 8 = 24\text{ cm}^2$.
		\end{itemchoice}
	}
\end{ex}

\begin{bt}%[9H5V17-1]
% nguon: SGK Toán 9 KNTT tập 1 – Bài 17, Vị trí tương đối của hai đường tròn
	Cho hai đường tròn $(O; 13\text{ cm})$ và $(O'; 15\text{ cm})$ cắt nhau tại $A, B$ sao cho dây chung $AB = 24\text{ cm}$ và hai tâm $O, O'$ nằm về hai phía đối với dây $AB$. Tính khoảng cách $OO'$ giữa hai tâm (đơn vị cm).
	\shortans{14}
	\loigiai{
		Gọi $H = OO' \cap AB$, ta có $OO' \perp AB$ tại $H$ và $AH = \dfrac{AB}{2} = 12\text{ cm}$.\\
		Trong tam giác $OHA$ vuông tại $H$: $OH = \sqrt{13^2 - 12^2} = 5\text{ cm}$.\\
		Trong tam giác $O'HA$ vuông tại $H$: $O'H = \sqrt{15^2 - 12^2} = 9\text{ cm}$.\\
		Vì $O$ và $O'$ nằm về hai phía của dây $AB$ nên $OO' = OH + O'H = 5 + 9 = 14\text{ cm}$.
	}
\end{bt}

\begin{bt}%[9H5V17-1]
% nguon: SGK Toán 9 KNTT tập 1 – Bài 17, Vị trí tương đối của hai đường tròn
	Cho hai đường tròn $(O; 13\text{ cm})$ và $(O'; 15\text{ cm})$ cắt nhau tại $A, B$ sao cho dây chung $AB = 24\text{ cm}$ và hai tâm $O, O'$ nằm về cùng một phía đối với dây $AB$. Tính khoảng cách $OO'$ giữa hai tâm (đơn vị cm).
	\shortans{4}
	\loigiai{
		Gọi $H$ là giao điểm của đường thẳng $OO'$ với dây $AB$, ta có $AH = \dfrac{AB}{2} = 12\text{ cm}$.\\
		Tính được $OH = \sqrt{13^2 - 12^2} = 5\text{ cm}$ và $O'H = \sqrt{15^2 - 12^2} = 9\text{ cm}$.\\
		Vì $O$ và $O'$ nằm cùng một phía đối với $AB$ nên $OO' = |O'H - OH| = 9 - 5 = 4\text{ cm}$.
	}
\end{bt}

% --- 9H5?17-2: Tiếp tuyến chung, bài toán tổng hợp ---

\begin{ex}%[9H5N17-2]
% nguon: SGK Toán 9 KNTT tập 1 – Bài 17, Vị trí tương đối của hai đường tròn
	Hai đường tròn tiếp xúc ngoài có tất cả bao nhiêu tiếp tuyến chung?
	\choice
	{$1$}
	{$2$}
	{\True $3$}
	{$4$}
	\loigiai{
		Hai đường tròn tiếp xúc ngoài có $2$ tiếp tuyến chung ngoài và $1$ tiếp tuyến chung trong đi qua tiếp điểm, tổng cộng là $3$ tiếp tuyến chung.
	}
\end{ex}

\begin{ex}%[9H5N17-2]
% nguon: SGK Toán 9 KNTT tập 1 – Bài 17, Vị trí tương đối của hai đường tròn
	Hai đường tròn cắt nhau có tất cả bao nhiêu tiếp tuyến chung?
	\choice
	{$0$}
	{$1$}
	{\True $2$}
	{$3$}
	\loigiai{
		Hai đường tròn cắt nhau chỉ có $2$ tiếp tuyến chung ngoài, không có tiếp tuyến chung trong.
	}
\end{ex}

\begin{ex}%[9H5H17-2]
% nguon: SGK Toán 9 KNTT tập 1 – Bài 17, Vị trí tương đối của hai đường tròn
	Cho hai đường tròn $(O; 16\text{ cm})$ và $(O'; 9\text{ cm})$ tiếp xúc ngoài tại điểm $A$. Gọi $DE$ là một tiếp tuyến chung ngoài của hai đường tròn ($D \in (O)$, $E \in (O')$). Xét tính đúng, sai của các khẳng định sau:
	\choiceTF
	{\True Độ dài đoạn nối tâm $OO'$ bằng $25\text{ cm}$}
	{\True Đường thẳng đi qua $A$ và vuông góc với $OO'$ là tiếp tuyến chung trong của hai đường tròn}
	{Độ dài đoạn tiếp tuyến chung ngoài $DE$ bằng $20\text{ cm}$}
	{\True Diện tích hình thang vuông $ODEO'$ bằng $300\text{ cm}^2$}
	\loigiai{
		\begin{itemchoice}
			\itemch Đúng. Vì hai đường tròn tiếp xúc ngoài nên $OO' = 16 + 9 = 25\text{ cm}$.
			\itemch Đúng. Đường thẳng qua $A$ vuông góc với $OO'$ tại $A$ vuông góc với cả bán kính $OA$ và $O'A$ nên là tiếp tuyến chung trong của hai đường tròn.
			\itemch Sai. Kẻ $O'H \perp OD$ tại $H$, ta có $OH = 16 - 9 = 7\text{ cm}$ nên $DE = O'H = \sqrt{25^2 - 7^2} = \sqrt{576} = 24\text{ cm}$.
			\itemch Đúng. Diện tích hình thang vuông $ODEO'$ là $\dfrac{16 + 9}{2} \cdot 24 = 300\text{ cm}^2$.
		\end{itemchoice}
	}
\end{ex}

\begin{ex}%[9H5H17-2]
% nguon: SGK Toán 9 KNTT tập 1 – Bài 17, Vị trí tương đối của hai đường tròn
	Cho hai đường tròn $(O; 7\text{ cm})$ và $(O'; 3\text{ cm})$ có khoảng cách hai tâm $OO' = 15\text{ cm}$. Xét tính đúng, sai của các khẳng định sau:
	\choiceTF
	{\True Hai đường tròn $(O)$ và $(O')$ ở ngoài nhau}
	{\True Hai đường tròn có $2$ tiếp tuyến chung ngoài và $2$ tiếp tuyến chung trong}
	{Khoảng cách ngắn nhất giữa một điểm trên $(O)$ và một điểm trên $(O')$ bằng $4\text{ cm}$}
	{\True Độ dài đoạn tiếp tuyến chung trong nối hai tiếp điểm của hai đường tròn bằng $5\sqrt{5}\text{ cm}$}
	\loigiai{
		\begin{itemchoice}
			\itemch Đúng. Vì $OO' = 15 > 7 + 3 = 10$ nên hai đường tròn ở ngoài nhau.
			\itemch Đúng. Hai đường tròn ở ngoài nhau luôn có $2$ tiếp tuyến chung ngoài và $2$ tiếp tuyến chung trong.
			\itemch Sai. Khoảng cách ngắn nhất giữa hai điểm trên hai đường tròn là $OO' - R - r = 15 - 7 - 3 = 5\text{ cm}$.
			\itemch Đúng. Độ dài đoạn tiếp tuyến chung trong là $\sqrt{OO'^2 - (R + r)^2} = \sqrt{15^2 - 10^2} = \sqrt{125} = 5\sqrt{5}\text{ cm}$.
		\end{itemchoice}
	}
\end{ex}

\begin{bt}%[9H5V17-2]
% nguon: SGK Toán 9 KNTT tập 1 – Bài 17, Vị trí tương đối của hai đường tròn
	Cho hai đường tròn $(O; 10\text{ cm})$ và $(O'; 2\text{ cm})$ ở ngoài nhau có khoảng cách hai tâm $OO' = 17\text{ cm}$. Tính độ dài đoạn tiếp tuyến chung ngoài nối hai tiếp điểm của hai đường tròn đó (đơn vị cm).
	\shortans{15}
	\loigiai{
		Gọi $A \in (O)$ và $B \in (O')$ là hai tiếp điểm của tiếp tuyến chung ngoài. Kẻ $O'H \perp OA$ tại $H$, khi đó $HABO'$ là hình chữ nhật nên $AB = O'H$ và $OH = R - r = 10 - 2 = 8\text{ cm}$.\\
		Trong tam giác $OHO'$ vuông tại $H$: $AB = O'H = \sqrt{OO'^2 - OH^2} = \sqrt{17^2 - 8^2} = \sqrt{225} = 15\text{ cm}$.
	}
\end{bt}

\begin{bt}%[9H5V17-2]
% nguon: SGK Toán 9 KNTT tập 1 – Bài 17, Vị trí tương đối của hai đường tròn
	Cho hai đường tròn $(O; 6\text{ cm})$ và $(O'; 3\text{ cm})$ ở ngoài nhau có khoảng cách hai tâm $OO' = 15\text{ cm}$. Tính độ dài đoạn tiếp tuyến chung trong nối hai tiếp điểm của hai đường tròn đó (đơn vị cm).
	\shortans{12}
	\loigiai{
		Gọi $A \in (O)$ và $B \in (O')$ là hai tiếp điểm của tiếp tuyến chung trong. Kẻ $O'K \perp OA$ tại $K$ (trên tia đối của tia $AO$), ta có $OK = R + r = 6 + 3 = 9\text{ cm}$ và $AB = O'K$.\\
		Trong tam giác $OKO'$ vuông tại $K$: $AB = O'K = \sqrt{OO'^2 - OK^2} = \sqrt{15^2 - 9^2} = \sqrt{144} = 12\text{ cm}$.
	}
\end{bt}
